\section{Method}
\label{sec:method}

% --- 1.5 pages ---

\subsection{Classical HRP recap}
\label{sec:hrp_recap}

Hierarchical Risk Parity \citep{LopezdePrado2016} operates in three stages.
Given an annualized covariance $\Sigma \in \mathbb{R}^{n \times n}$, it
(i) computes a correlation distance
$D_{ij} = \sqrt{(1 - \rho_{ij})/2}$ where
$\rho_{ij} = \Sigma_{ij}/\sqrt{\Sigma_{ii}\Sigma_{jj}}$,
(ii) builds a binary tree by single-linkage agglomerative clustering on $D$,
and (iii) recursively bisects the tree, allocating capital between the two
sub-trees in inverse proportion to their cluster volatilities.
Each step is non-differentiable: the agglomerative clustering produces a
combinatorial tree, and the recursive bisection involves a hard
volatility-weighted split. Consequently HRP cannot be jointly trained
with upstream feature extractors, regardless of the loss function on top.
DHRP replaces both stages with smooth parameterized counterparts that
preserve the hierarchical inductive bias.

\subsection{DHRP: a differentiable soft-gating tree}
\label{sec:dhrp_layer}

DHRP maps a feature vector $x_t \in \mathbb{R}^f$ and an annualized
covariance matrix $\Sigma_t \in \mathbb{R}^{n \times n}$ to portfolio
weights $w_t \in \Delta^{n-1}$ via four learnable components:
% TODO: figure 1 - architecture diagram
% \begin{figure}[t]
%   \centering
%   \includegraphics[width=0.95\textwidth]{architecture.pdf}
%   \caption{DHRP and LLM-DHRP architectures.}
%   \label{fig:arch}
% \end{figure}

\paragraph{Feature fusion.}
A cov-projection MLP $f_{\text{cov}}: \mathbb{R}^{n^2} \to \mathbb{R}^f$
encodes the flattened normalized covariance, then layer-norm fuses with
price features:
\begin{equation}
h_t = \mathrm{LayerNorm}\left(x_t + f_{\text{cov}}\!\left(
\frac{\mathrm{vec}(\Sigma_t)}{\|\Sigma_t\|_\infty + \epsilon}\right)\right).
\end{equation}

\paragraph{Soft-gating tree.}
A perfect binary tree of depth $L$ has $2^L$ leaves and $2^L-1$ internal
nodes. Each internal node $k$ has a gate
$g_k: \mathbb{R}^f \to \mathbb{R}^2$ producing routing probabilities:
\begin{equation}
\pi_k = \mathrm{softmax}(g_k(h_t) / \tau_k),
\quad \tau_k = \exp(\theta_{\tau,k}) \in (0.1, 5.0].
\label{eq:gates}
\end{equation}
The leaf budget for leaf $\ell$ along path $P_\ell$ is
$p_\ell = \prod_{(k, b) \in P_\ell} \pi_k[b]$, then normalized by a
learned leaf prior $\beta = \mathrm{softmax}(\theta_\beta)$:
$b_\ell = \beta_\ell p_\ell / \sum_j \beta_j p_j$.

\paragraph{Soft asset assignment.}
A learnable matrix $A \in \mathbb{R}^{n \times 2^L}$ row-normalizes via
softmax; asset $i$'s share of leaf budgets is
\begin{equation}
a_i = \sum_{\ell} \mathrm{softmax}(A_{i,:})_\ell \cdot b_\ell.
\end{equation}

\paragraph{Volatility blending.}
Final weights blend the tree-derived $a$ with inverse-volatility:
\begin{equation}
w_i = \alpha \cdot a_i + (1-\alpha) \cdot
\frac{1/\sigma_i}{\sum_j 1/\sigma_j},
\qquad \alpha = \sigma(\theta_\alpha),
\label{eq:blend}
\end{equation}
followed by a non-negativity clip and re-normalization.

\paragraph{Loss.}
A multi-objective loss combines CRRA utility, smooth Sharpe, HRP
regularization (curriculum-scheduled), risk penalty, concentration (HHI),
and turnover penalty:
\begin{equation}
\mathcal{L}_{\text{DHRP}} = -(\mathrm{CRRA} + \mathrm{Sharpe} +
\lambda_e \mathrm{Entropy}) +
\lambda_h \|w - w_{\text{HRP}}\|_2^2 +
\lambda_r w^\top \Sigma w +
\lambda_c \mathrm{HHI}(w) +
\lambda_t \mathrm{TO}(w).
\end{equation}
The HRP regularizer anchors $w$ to classical HRP weights $w_{\text{HRP}}$
with $\lambda_h$ decayed from $0.3 \to 0.05$ over training, providing a
warm-start that recovers HRP at $\tau \to 0$ (Proposition~\ref{prop:recovery}).

\begin{proposition}[HRP recovery]
\label{prop:recovery}
For any covariance $\Sigma$ such that the agglomerative clustering yields
a unique tree $T^*$, there exists a parameter setting
$(\theta_\tau, \theta_\beta, A) \to \theta^*$ and temperature
$\tau_k \to 0$ under which the DHRP layer (Eqs.~\ref{eq:gates}--\ref{eq:blend})
produces weights identical to classical HRP applied to $\Sigma$.
\end{proposition}

\noindent\textit{Proof sketch.} Hard-route gates concentrate on the
HRP cluster split direction; the soft-assignment $A$ with one-hot rows
recovers leaf membership; $\alpha = 1$ disables the inverse-vol blend.
Full proof in Appendix~\ref{app:proof}.

\subsection{LLM-DHRP: cross-modal fusion}
\label{sec:llm_dhrp_arch}

LLM-DHRP extends DHRP by injecting per-asset, time-varying text features
$T_t \in \mathbb{R}^{n \times d_T}$ derived from financial news headlines.
Headlines for asset $a$ within a 90-day rolling window are encoded by
FinBERT \citep{Huang2023} ($d_T = 768$), L2-normalized per asset, then
aggregated via concatenated mean and max-pool to $\mathbb{R}^{2 d_T}$
(see Section~\ref{sec:benchmark} for coverage details). A
TextProjection MLP maps to $\mathbb{R}^f$, fused with $h_t$ via gated
cross-attention:
\begin{equation}
\hat h_t = h_t + g_{\text{text}} \cdot
\mathrm{CrossAttn}(h_t, T_t^{\text{proj}}, T_t^{\text{proj}}),
\quad g_{\text{text}} = \sigma(\theta_g - 1.0),
\end{equation}
with gate bias $-1.0$ initially suppressing text. Modality dropout
(probability 0.1) randomly masks text during training, preventing
pathway collapse. Training uses warm-start from a pretrained DHRP:
phase~1 freezes price-pathway parameters and trains text/fusion only;
phase~2 unfreezes all with reduced learning rate. The full training
recipe and ablations on fusion variants are in
Appendix~\ref{app:hyperparams}.
